% ---- Opci zakljucak (ne numerise se) ---------------------------------------
\chapter*{Zaklju\v{c}ak}
\addcontentsline{toc}{chapter}{Zaklju\v{c}ak}

U radu je prikazano portanje aplikacije TetriSaraj sa golog izvr\v{s}avanja na
operativni sistem za rad u realnom vremenu FreeRTOS, na otvorenoj RISC-V
platformi realizovanoj kao meki procesor PicoRV32 na plo\v{c}i Tang Nano 9K. Na
osnovu izlo\v{z}enog mo\v{z}e se zaklju\v{c}iti da je sva tri pitanja postavljena
u uvodu mogu\'ce potvrdno rije\v{s}iti. Aplikacija je uspje\v{s}no razdvojena na
\v{c}iste zadatke pod upravom raspore\dj{}iva\v{c}a, a njen rad je dokazan na
stvarnom hardveru putem serijskog ispisa u kojem logika igre napreduje. Cijeli
sistem, zajedno sa video izlazom, uklopljen je u stroge resurse plo\v{c}e i
pokrenut na hardveru. Kona\v{c}no, izmjereno je vrijeme zamjene konteksta na
stvarnom hardveru i oprezno upore\dj{}eno sa rezultatima iz simulatora.

Posebno je istaknuto otkri\'ce da uzrok neuspjelog uklapanja nije bila
koli\v{c}ina logike, nego nedostatak blokovske memorije, kao i to da alat za
sintezu pouzdano koristi blokovsku memoriju samo za dubine koje su stepeni broja
dva. Ta dva uvida bila su klju\v{c}na za uspje\v{s}no uklapanje sistema. Tako\dj{}er
je pokazano da su \v{c}iste vrijednosti zamjene konteksta na hardveru i u
simulatoru istog reda veli\v{c}ine.

Iako rad obuhvata cijeli put od porta operativnog sistema do mjerenja na
hardveru, pojedina pitanja ostaju otvorena za daljnje istra\v{z}ivanje.
Vizuelna potvrda kona\v{c}ne slike igre na ekranu nije zatvorena u okviru ovog
rada. Potpuno interaktivno upravljanje ograni\v{c}eno je malim brojem fizi\v{c}kih
dugmadi na plo\v{c}i. Tako\dj{}er, mjerenja bi se mogla pro\v{s}iriti na tro\v{s}ak
samog takta raspore\dj{}iva\v{c}a i na pona\v{s}anje sistema pod ve\'cim
optere\'cenjem. Ovi pravci predstavljaju prirodan nastavak rada.
